\section*{Chapter \ref{ch:g10:ecosystems} --- Ecosystems}

\begin{solution}{exo:g10:ecosystems:1}
In a forest, a \omterm{def:g10:ecosystems:levels}{population} might be all oaks of one species in the stand. The
\omterm{def:g10:ecosystems:levels}{community} is oaks plus other plants, animals, fungi and microbes interacting
there. The \omterm{def:g10:ecosystems:levels}{ecosystem} is that \omterm{def:g10:ecosystems:levels}{community} together with abiotic factors: soil,
climate, water, light and nutrients.
\end{solution}

\begin{solution}{exo:g10:ecosystems:2}
A \omterm{def:g10:ecosystems:trophic-roles}{producer} fixes energy into organic matter (e.g.\ grass by \omterm{def:g10:photosynthesis:def}{photosynthesis}).
A \omterm{def:g10:ecosystems:trophic-roles}{primary consumer} eats \omterm{def:g10:ecosystems:trophic-roles}{producers} (e.g.\ grasshopper or rabbit). A \omterm{def:g10:ecosystems:trophic-roles}{decomposer}
breaks down dead material and recycles nutrients (e.g.\ soil bacterium or
fungus on litter).
\end{solution}

\begin{solution}{exo:g10:ecosystems:3}
Redraw as seeds $\to$ mouse $\to$ snake $\to$ hawk, with arrows pointing toward
the \omterm{def:g10:ecosystems:trophic-roles}{consumer} (direction of energy flow). \omterm{def:g10:ecosystems:trophic}{Trophic levels}: seeds/\omterm{def:g10:ecosystems:trophic-roles}{producer} (1),
mouse (2), snake (3), hawk (4).
\end{solution}

\begin{solution}{exo:g10:ecosystems:4}
Only a fraction of energy at one \omterm{def:g10:ecosystems:trophic}{trophic level} becomes biomass at the next;
much is lost as heat in respiration and as unassimilated material. Too little
energy remains to support arbitrarily long chains, so \omterm{def:g10:ecosystems:food-web}{food chains} are usually
short.
\end{solution}

\begin{solution}{exo:g10:ecosystems:5}
Processes that remove atmospheric $\mathrm{CO_2}$ include \omterm{def:g10:photosynthesis:def}{photosynthesis} (and,
at broader scale, ocean uptake). Processes that add $\mathrm{CO_2}$ include
respiration, decomposition and combustion of biomass or fossil fuels.
\end{solution}

\begin{solution}{exo:g10:ecosystems:6}
If about $10\%$ passes each step, \omterm{def:g10:ecosystems:trophic-roles}{primary consumers} receive roughly
$\qty{1000}{kJ}$ from $\qty{10000}{kJ}$ at \omterm{def:g10:ecosystems:trophic-roles}{producers}, and \omterm{def:g10:ecosystems:trophic-roles}{secondary consumers}
receive about $\qty{100}{kJ}$ after a second transfer.
\end{solution}

\begin{solution}{exo:g10:ecosystems:7}
\omterm{def:g10:ecosystems:nitrogen}{Nitrogen-fixing} microbes convert atmospheric $\mathrm{N_2}$ into forms plants
can assimilate. Legumes host fixers in root nodules; rotating legumes into a
cropping system adds usable nitrogen to soil and can reduce fertiliser need
for following crops.
\end{solution}

\begin{solution}{exo:g10:ecosystems:8}
Insect-eating fish lose a major food source and may decline; birds that eat
those fish may decline in turn. Surviving insects might include resistant
\omterm{def:g10:ecosystems:levels}{populations}; some fish might switch prey. If the pesticide persists and
biomagnifies, top predators can accumulate higher doses.
\end{solution}

\begin{solution}{exo:g10:ecosystems:9}
\omterm{def:g10:ecosystems:habitat-niche}{Habitat} is the place --- e.g.\ intertidal rock in the splash or tidal zone.
\omterm{def:g10:ecosystems:habitat-niche}{Niche} is the barnacle's role and resource use: filter feeding at certain tidal
heights, competing for space, and tolerating desiccation and salinity ranges.
\end{solution}

\begin{solution}{exo:g10:ecosystems:10}
\omterm{def:g10:scientific-biology:variables}{Independent variable}: soil moisture (wet vs dry). \omterm{def:g10:scientific-biology:variables}{Dependent variable}: mass
loss of litter over a fixed time (or another decomposition measure). \omterm{def:g10:scientific-biology:variables}{Controlled
variables}: temperature, litter type, mesh size, starting mass, duration.
Expect faster decomposition when moist enough for microbial activity, unless
waterlogging limits oxygen.
\end{solution}

\begin{solution}{exo:g10:ecosystems:11}
Example path: \omterm{def:g10:ecosystems:nitrogen}{nitrogen fixation} $\to$ plant uptake of ammonium or nitrate $\to$
plant \omterm{def:g10:molecules-of-life:proteins}{protein} $\to$ cow eats plant $\to$ cow muscle \omterm{def:g10:molecules-of-life:proteins}{protein}. Return to the
atmosphere can occur by \omterm{def:g10:ecosystems:nitrogen}{denitrification} after excretion and decomposition,
releasing $\mathrm{N_2}$ (or $\mathrm{N_2O}$).
\end{solution}

\begin{solution}{exo:g10:ecosystems:12}
A \omterm{def:g10:ecosystems:keystone}{keystone species} has a \omterm{def:g10:ecosystems:levels}{community} effect out of proportion to its abundance.
Sea otters may be few, but by eating urchins they protect kelp forests that
house many species; without otters, urchins can destroy kelp \omterm{def:g10:ecosystems:habitat-niche}{habitat}. Counting
individuals would miss that structural role.
\end{solution}

\begin{solution}{exo:g10:ecosystems:13}
\omterm{def:g10:ecosystems:succession}{Primary succession} starts on new \omterm{def:g10:membranes-enzymes:enzyme}{substrate} without soil (bare rock, lava);
\omterm{def:g10:ecosystems:succession}{secondary succession} follows disturbance that leaves soil (fire, abandoned
field). After fire: herbs and grasses often colonise first, then shrubs, then
trees if climate allows --- a \omterm{def:g10:ecosystems:trophic-roles}{secondary} \omterm{def:g10:molecules-of-life:proteins}{sequence} on remaining soil. Patchy
disturbance creates a mosaic of stages, which can raise landscape \omterm{def:g10:biodiversity-classification:biodiversity}{species
diversity} compared with a single uniform stage.
\end{solution}

\begin{solution}{pb:g10:ecosystems:1}
Answers depend on site; full credit requires:
\textbf{Part I.} At least eight groups with trophic roles.
\textbf{Part II.} A web with at least six arrows; the longest chain identified
with transfer count; an energy-loss argument that predator biomass is less
than prey biomass over the long term.
\textbf{Part III.} One plausible human effect on carbon or nitrogen; without
\omterm{def:g10:ecosystems:trophic-roles}{decomposers}, dead matter accumulates and nutrient recycling stalls; one
conservation action linked to \omterm{def:g10:ecosystems:habitat-niche}{habitat}, resources or connectivity.
\end{solution}
